% !TeX encoding = UTF-8
\documentclass[variant=docente,cover=institutional,mode=screen]{../cls/ua-engineering-v6-article}
% !TeX encoding = UTF-8
\UASetUniversity{Universidad Austral}
\UASetFaculty{Facultad de Ingeniería}
\UASetDepartment{Departamento de Computación}
\UASetCampus{Pilar}
\UASetCareer{Ingeniería Informática}
\UASetCommission{Comisión A}
\UASetProfessor{Equipo docente}
\UASetSemester{Segundo cuatrimestre 2026}
\UASetCourse{Sistemas Distribuidos}
\UASetDate{2026}

\UASetClassNumber{05}
\UASetClassTitle{Tiempo y causalidad}
\title{Clase 05 --- Tiempo y causalidad}
\author{\UAProfessor}
\date{\UADate}
\begin{document}
\maketitle

\section{Objetivo pedagógico profundo}
Que el estudiante reemplace la intuición de reloj global por la noción de dependencia causal. El objetivo no es memorizar Lamport y vector clocks, sino saber qué permiten inferir sobre orden, concurrencia y conflictos.

\section{Resultados de aprendizaje esperados}
\begin{itemize}
\item explicar por qué el tiempo físico no alcanza para ordenar eventos distribuidos;
\item aplicar la relación happened-before;
\item calcular relojes de Lamport en una ejecución simple;
\item explicar el límite de Lamport para detectar concurrencia;
\item comparar vector clocks y decidir cuándo dos eventos son concurrentes.
\end{itemize}

\section{Guion sugerido}
\begin{enumerate}
\item Abrir con dos eventos en nodos distintos y preguntar cuál ocurrió primero.
\item Mostrar que timestamps físicos no prueban causalidad.
\item Definir happened-before con orden local, mensajes y transitividad.
\item Calcular Lamport y discutir su límite.
\item Introducir vector clocks como mecanismo para detectar concurrencia.
\item Cerrar con conflictos concurrentes en stores eventualmente consistentes.
\end{enumerate}

\section{Errores conceptuales frecuentes}
\begin{itemize}
\item creer que NTP resuelve orden causal;
\item confundir orden total arbitrario con causalidad;
\item inferir causalidad solo porque un timestamp Lamport es menor;
\item olvidar que vector clocks tienen costo de metadata.
\end{itemize}

\section{Preguntas disparadoras}
\begin{itemize}
\item ¿Existe un mensaje o camino causal entre estos eventos?
\item ¿Qué puede concluir Lamport y qué no?
\item ¿Son comparables estos vectores?
\item ¿Qué debería hacer una aplicación ante escrituras concurrentes?
\end{itemize}

\begin{uawarn}
Criterio de logro docente: la clase salió bien si el estudiante puede distinguir ``ocurrió antes en el reloj'' de ``causó o precedió lógicamente''.
\end{uawarn}
\end{document}
